\chapter{Ngẩng lên nào}
\begin{tinhhuongbox}{Tình huống | Situation}
\songngu{Lớp vừa vào tiết nhưng mắt nhiều em còn ở ngoài hành lang.}{The class has just started, but many eyes are still in the hallway.}
\songngu{Nếu thầy cô nói luôn một đoạn dài, lời nói sẽ đi qua lớp như gió đi qua cửa sổ.}{If you start with a long speech right away, your words may pass through the class like wind through a window.}
\end{tinhhuongbox}
\begin{noitrenlop}
\songngu{Ngẩng lên nào, mình bắt đầu gọn thôi.}{Eyes up, everyone. We will start simply.}
\end{noitrenlop}
\begin{vihieuqua}
\songngu{Câu này ngắn, rõ và không có mùi càm ràm.}{This line is short, clear, and does not sound like nagging.}
\end{vihieuqua}
\begin{englishpocket}
\songngu{Cụm nên nhớ: \textit{Eyes up, everyone.}}{Phrase to keep: \textit{Eyes up, everyone.}}
\end{englishpocket}
\chapter{Thân vào lớp rồi, não vào chưa?}
\begin{tinhhuongbox}{Tình huống | Situation}
\songngu{Sau giờ ra chơi, nhiều em ngồi đúng chỗ nhưng đầu óc vẫn còn treo ngoài sân.}{After recess, many students sit in the right place, but their minds are still outside.}
\songngu{Lúc này một câu hơi vui thường hiệu quả hơn một câu hơi gắt.}{At this point, a slightly funny line often works better than a slightly harsh one.}
\end{tinhhuongbox}
\begin{noitrenlop}
\songngu{Thân vào lớp rồi, não vào chưa, hay tôi phải đi đón nốt?}{Your body is here, but is your brain here too, or do I need to go pick it up?}
\end{noitrenlop}
\begin{vihieuqua}
\songngu{Lớp thường bật cười nhẹ và tự hiểu mình cần quay về tiết học ngay.}{The class usually smiles a little and instantly gets the message to return to the lesson.}
\end{vihieuqua}
\begin{englishpocket}
\songngu{Cụm nên nhớ: \textit{Do I need to go pick it up?}}{Phrase to keep: \textit{Do I need to go pick it up?}}
\end{englishpocket}
\chapter{Chọn mood để vào tiết}
\begin{tinhhuongbox}{Tình huống | Situation}
\songngu{Có những tiết lớp không hẳn ồn, chỉ là hơi nguội và hơi lười khởi động.}{Some lessons are not noisy; they are just a little cold and a little slow to start.}
\songngu{Thầy cô cần một câu đủ nhẹ để cả lớp nhập cuộc cùng lúc.}{You need a light enough prompt for the whole class to join at once.}
\end{tinhhuongbox}
\begin{noitrenlop}
\songngu{Hôm nay vào tiết theo mood nào: tỉnh táo, tử tế, hay cứu gấp?}{Which mood are we using today: sharp, kind, or emergency rescue?}
\end{noitrenlop}
\begin{vihieuqua}
\songngu{Câu hỏi vui nhưng vẫn buộc lớp tự chọn thái độ cho 5 phút đầu.}{The question is playful, but it still makes the class choose an attitude for the first five minutes.}
\end{vihieuqua}
\begin{englishpocket}
\songngu{Cụm nên nhớ: \textit{Which mood are we using today?}}{Phrase to keep: \textit{Which mood are we using today?}}
\end{englishpocket}
\chapter{Đổi vai hai mươi giây}
\begin{tinhhuongbox}{Tình huống | Situation}
\songngu{Khi lớp nghe quá thụ động, tốt nhất là trả lại mic cho học sinh trong vài giây.}{When the class is listening too passively, the best move is to give the microphone back to students for a few seconds.}
\songngu{Đổi vai ngắn làm không khí bật lên rất nhanh.}{A short role switch can lift the room very quickly.}
\end{tinhhuongbox}
\begin{noitrenlop}
\songngu{Hai mươi giây thôi: nếu em là thầy cô, em sẽ mở tiết này bằng câu nào?}{Just twenty seconds: if you were the teacher, what line would you use to open this lesson?}
\end{noitrenlop}
\begin{vihieuqua}
\songngu{Học sinh vừa cười vừa phải nghĩ, và đó là lúc lớp thật sự vào bài.}{Students laugh and think at the same time, and that is when the class truly enters the lesson.}
\end{vihieuqua}
\begin{englishpocket}
\songngu{Cụm nên nhớ: \textit{What line would you use?}}{Phrase to keep: \textit{What line would you use?}}
\end{englishpocket}
\chapter{Ba từ cho bài hôm nay}
\begin{tinhhuongbox}{Tình huống | Situation}
\songngu{Đầu giờ dài quá dễ làm lớp trôi ngay từ bậc cửa.}{An overlong opening can make the class drift away right at the doorway.}
\songngu{Một nhiệm vụ nhỏ sẽ kéo sự chú ý về rất nhanh.}{A tiny task can pull attention back very fast.}
\end{tinhhuongbox}
\begin{noitrenlop}
\songngu{Mỗi bàn cho tôi ba từ dự đoán về bài hôm nay, ngắn thôi, đừng viết tiểu thuyết.}{Each desk gives me three words to predict today's lesson, keep it short, do not write a novel.}
\end{noitrenlop}
\begin{vihieuqua}
\songngu{Nói ba từ thì ai cũng dám, mà ba từ lại đủ để lớp bắt đầu nghĩ.}{Anyone dares to say three words, and three words are enough to make the class start thinking.}
\end{vihieuqua}
\begin{englishpocket}
\songngu{Cụm nên nhớ: \textit{Keep it short.}}{Phrase to keep: \textit{Keep it short.}}
\end{englishpocket}